% ================================================================
%  Rapport — Application Traveling (TravelShare + TravelPath)
%  Compilateur : pdfLaTeX — Overleaf
% ================================================================
\documentclass[a4paper,12pt,openany]{report}

% ── Encodage & langue ──────────────────────────────────────────
\usepackage[utf8]{inputenc}
\usepackage[T1]{fontenc}
\usepackage[french]{babel}
\usepackage{lmodern}
\usepackage{inconsolata}

% ── Mise en page ───────────────────────────────────────────────
\usepackage[top=2.5cm, bottom=2.5cm, left=2.8cm, right=2.8cm]{geometry}
\usepackage{parskip}
\usepackage{microtype}
\setlength{\parindent}{0pt}
\setlength{\parskip}{6pt}

% ── Couleurs (palette sobre) ────────────────────────────────────
\usepackage[dvipsnames,table]{xcolor}
\definecolor{navy}{HTML}{2C4A6E}
\definecolor{teal}{HTML}{2C7873}
\definecolor{terra}{HTML}{B85C38}
\definecolor{codebg}{HTML}{F3F4F6}
\definecolor{codetext}{HTML}{1F2937}
\definecolor{lightbg}{HTML}{EEF2F7}
\definecolor{greentxt}{HTML}{276749}
\definecolor{warnbg}{HTML}{FEFCE8}
\definecolor{okbg}{HTML}{F0FDF4}
\definecolor{infobg}{HTML}{EFF6FF}
\definecolor{errbg}{HTML}{FEF2F2}
\definecolor{rowalt}{HTML}{F9FAFB}
\definecolor{bord}{HTML}{D1D5DB}

% ── Titres (sobres et lisibles) ─────────────────────────────────
\usepackage{titlesec}
\titleformat{\chapter}[block]
  {\normalfont\Large\bfseries\color{navy}}
  {\thechapter.\enspace}{0pt}{}
  [\vspace{3pt}{\color{bord}\rule{\textwidth}{0.8pt}}\vspace{2pt}]
\titlespacing*{\chapter}{0pt}{18pt}{12pt}

\titleformat{\section}
  {\normalfont\large\bfseries\color{navy}}
  {\thesection\enspace}{0pt}{}
\titlespacing*{\section}{0pt}{14pt}{4pt}

\titleformat{\subsection}
  {\normalfont\normalsize\bfseries\color{black}}
  {\thesubsection\enspace}{0pt}{}
\titlespacing*{\subsection}{0pt}{10pt}{2pt}

% ── En-têtes ───────────────────────────────────────────────────
\usepackage{fancyhdr}
\pagestyle{fancy}\fancyhf{}
\fancyhead[L]{\small\color{gray!80}Traveling — TravelShare \& TravelPath}
\fancyhead[R]{\small\color{gray!80}Hammar · Soufyani — M1 2025-2026}
\fancyfoot[C]{\small\color{gray}\thepage}
\renewcommand{\headrulewidth}{0.3pt}

% ── Liens ──────────────────────────────────────────────────────
\usepackage[colorlinks=true, linkcolor=navy, urlcolor=teal]{hyperref}

% ── Tableaux ───────────────────────────────────────────────────
\usepackage{booktabs}
\usepackage{array}
\usepackage{tabularx}
\usepackage{longtable}
\renewcommand{\arraystretch}{1.35}

% ── Code source (fond clair, lisible) ──────────────────────────
\usepackage{listings}
\lstdefinestyle{java}{
  backgroundcolor  = \color{codebg},
  basicstyle       = \ttfamily\footnotesize\color{codetext},
  keywordstyle     = \color{navy}\bfseries,
  stringstyle      = \color{teal!80!black},
  commentstyle     = \color{gray}\itshape,
  numberstyle      = \tiny\color{gray!60},
  language         = Java,
  breaklines       = true,
  showstringspaces = false,
  frame            = single, framerule = 0.4pt,
  rulecolor        = \color{bord},
  xleftmargin = 8pt, xrightmargin = 6pt,
  aboveskip = 8pt, belowskip = 6pt,
  tabsize = 4, numbers = left, numbersep = 10pt,
  morekeywords = {LiveData,MediatorLiveData,AndroidViewModel,
    AppDatabase,ExecutorService,SensorManager,MediaRecorder,
    MediaPlayer,ActivityResultLauncher,TravelPlan,PlanStep,
    GeoPoint,Polyline,BoundingBox,SharedViewModel,TravelPathViewModel},
}
\lstset{style=java}

% ── Boîtes simples ─────────────────────────────────────────────
\usepackage[most]{tcolorbox}
\tcbset{mybox/.style={
  enhanced, arc=3pt, boxrule=0.5pt,
  left=8pt, right=8pt, top=5pt, bottom=5pt,
  fonttitle=\bfseries\small}}

\newenvironment{infobox}[1][Note]{%
  \begin{tcolorbox}[mybox, colback=infobg, colframe=navy!30,
    title=\textcolor{navy}{#1}]\small\color{navy!80}}
  {\end{tcolorbox}}

\newenvironment{okbox}[1][]{%
  \begin{tcolorbox}[mybox, colback=okbg, colframe=teal!40,
    title=\textcolor{greentxt}{#1}]\small\color{greentxt!80!black}}
  {\end{tcolorbox}}

\newenvironment{warnbox}[1][Attention]{%
  \begin{tcolorbox}[mybox, colback=warnbg, colframe=Goldenrod!60,
    title=\textcolor{brown!80}{#1}]\small\color{brown!70!black}}
  {\end{tcolorbox}}

\newenvironment{errbox}[1][Problème]{%
  \begin{tcolorbox}[mybox, colback=errbg, colframe=terra!50,
    title=\textcolor{terra}{#1}]\small\color{terra!80!black}}
  {\end{tcolorbox}}

% ── Commandes utiles ───────────────────────────────────────────
\newcommand{\ic}[1]{\texttt{\small\colorbox{lightbg}{\color{navy!90} #1}}}
\newcommand{\cyes}{\textcolor{greentxt}{\bfseries\checkmark}}
\newcommand{\cno}{\textcolor{terra}{\bfseries$\times$}}

% ── Captures d'écran réelles ───────────────────────────────────
\usepackage{graphicx}
\newcommand{\screenshot}[2][0.38\textwidth]{%
  \begin{center}
  \vspace{4pt}
  \includegraphics[width=#1,keepaspectratio]{#2}
  \vspace{4pt}
  \end{center}
}

% ── Enumération ────────────────────────────────────────────────
\usepackage{enumitem}
\usepackage{multicol}
\setlist[itemize]{noitemsep, topsep=3pt, leftmargin=1.3em}
\setlist[enumerate]{noitemsep, topsep=3pt}

% ==============================================================
\begin{document}
\pagestyle{empty}

% ──────────────────────────────────────────────────────────────
%  PAGE DE GARDE
% ──────────────────────────────────────────────────────────────
\begin{titlepage}
\centering
\vspace*{2cm}

{\color{navy}\rule{\textwidth}{1.5pt}}
\vspace{1cm}

{\fontsize{42}{50}\selectfont\bfseries\color{navy} Traveling\par}
\vspace{0.3cm}
{\large\color{teal} TravelShare \& TravelPath\par}
\vspace{0.2cm}
{\normalsize\color{gray} Application Android — Java\par}

\vspace{1cm}
{\color{navy}\rule{\textwidth}{0.5pt}}

\vspace{2.5cm}

\begin{tabular}{ll}
  \textbf{Étudiants} & Sofiane Hammar \quad — \quad TravelShare + Passerelle \\[4pt]
                     & Anouar Soufyani \quad — \quad TravelPath \\[14pt]
  \textbf{Formation} & Master 1 Programmation Mobile \\[4pt]
  \textbf{Année}     & 2025-2026 \\[14pt]
  \textbf{Dépôt Git} & \texttt{github.com/Soof2/Travelshare}
\end{tabular}

\vfill
{\color{bord}\rule{\textwidth}{0.5pt}}\\[4pt]
{\small\color{gray} Rapport de projet — Module Programmation Mobile}
\end{titlepage}

\nopagecolor
\pagestyle{fancy}
\tableofcontents
\clearpage

% ──────────────────────────────────────────────────────────────
%  RÉSUMÉ
% ──────────────────────────────────────────────────────────────
\begin{tcolorbox}[mybox, colback=lightbg, colframe=bord,
  left=14pt, right=14pt, top=12pt, bottom=12pt,
  title=\large\bfseries\color{navy}Résumé]

Dans le cadre du module Programmation Mobile (Master 1, 2025-2026), nous avons
conçu et développé \textbf{Traveling}, une application Android complète dédiée
au voyage, réalisée en binôme en Java.

L'idée de départ était de proposer deux services complémentaires autour du voyage.
\textbf{TravelShare} permet de partager et de découvrir des photos de voyages :
on peut publier une photo depuis sa galerie ou son appareil photo, la géolocaliser,
y associer une note vocale, et la partager avec un groupe ou en public. Les autres
utilisateurs peuvent la retrouver grâce à une recherche textuelle ou vocale, par
catégorie, par auteur, par période, ou même en secouant simplement leur téléphone
pour obtenir des suggestions aléatoires.

\textbf{TravelPath} complète cette expérience en générant des itinéraires personnalisés.
L'utilisateur renseigne une ville, ses activités préférées, son budget et sa disponibilité,
et l'application lui propose jusqu'à trois parcours adaptés — économique, équilibré ou confort —
avec de vrais lieux issus d'OpenStreetMap, un tracé routier réel, la météo du jour
et les adresses de chaque étape.

Sur le plan technique, tout repose sur le pattern MVVM avec une base de données Room,
et uniquement des services gratuits et open-source (Nominatim, OSRM, Overpass, Open-Meteo,
ML Kit). Aucune clé API payante n'est utilisée.

\vspace{6pt}
\begin{tabular}{ll}
  \textbf{Mots-clés :} &
  Android, Java, MVVM, Room, OSMDroid, OSRM, Nominatim, Overpass, \\
  & Open-Meteo, ML Kit, MediaRecorder, SensorManager, LiveData
\end{tabular}
\end{tcolorbox}

\clearpage

% ==============================================================
\chapter{Introduction}
% ==============================================================

\section{Contexte du projet}

Ce rapport présente l'application \textbf{Traveling}, développée dans le cadre du module
\textbf{Programmation Mobile} en Master 1, année 2025-2026. Le projet a été réalisé en
binôme par \textbf{Sofiane Hammar} et \textbf{Anouar Soufyani}.

L'objectif était de concevoir une application Android complète, en mettant en pratique
les concepts vus en cours : base de données locale, architecture MVVM, capteurs matériels,
notifications, cartographie et communication entre composants. Le projet devait également
démontrer une capacité à réaliser une application cohérente avec plusieurs modules
interdépendants.

\section{Ce que fait l'application}

\textbf{Traveling} est organisée autour de trois parties qui fonctionnent ensemble :

\begin{itemize}
  \item \textbf{TravelShare} est le côté réseau social de l'application. On y publie des
    photos de voyages, on les explore par catégorie, lieu, auteur ou période, on rejoint
    des groupes privés pour partager des albums, et on reçoit des notifications quand
    un lieu ou un auteur qu'on suit publie une nouvelle photo.

  \item \textbf{TravelPath} est le générateur d'itinéraires. On entre une ville et ses
    préférences (activités souhaitées, budget, durée, niveau d'effort), et l'application
    propose jusqu'à trois parcours avec de vrais restaurants, musées et parcs issus
    d'OpenStreetMap, un itinéraire tracé sur une carte, la météo du jour et les adresses
    de chaque étape.

  \item \textbf{La passerelle} relie les deux modules. Depuis une photo TravelShare, on
    peut générer un itinéraire vers ce lieu en un clic. Et depuis un parcours TravelPath,
    on peut voir les photos TravelShare associées à chaque étape.
\end{itemize}

\section{Répartition du travail}

\begin{tabularx}{\textwidth}{>{\bfseries}l >{\bfseries}l X}
  \toprule
  Étudiant & Partie & Fonctionnalités développées \\
  \midrule
  Sofiane Hammar
    & TravelShare
    & Authentification, publication de photos, explorateur multi-filtres,
      pull-to-refresh + pagination, accéléromètre (shake), note vocale,
      ML Kit tags IA, photo de profil + bio, partage/modifier/supprimer posts,
      groupes, chat, badges, notifications, carte, passerelle TS$\to$TP \\[4pt]
  Anouar Soufyani
    & TravelPath
    & Formulaire de génération, vrais lieux Overpass API, algorithme de parcours
      (durée/effort/budget dynamiques), géocodage Nominatim + Haversine,
      adresses réelles des étapes, tracé GPS OSRM, météo Open-Meteo,
      export PDF, régénération, suppression parcours, passerelle TP$\to$TS \\[4pt]
  Les deux
    & Architecture
    & Base de données Room, ViewModels partagés, modèles de données, intégration
      générale \\
  \bottomrule
\end{tabularx}


\section{Fonctionnalités implémentées}

Le tableau suivant récapitule l'ensemble des fonctionnalités réalisées dans
l'application.

% ── TravelShare ──────────────────────────────────────────────────
\vspace{6pt}
\begin{tcolorbox}[enhanced, colback=infobg, colframe=teal, arc=4pt,
  boxrule=0.6pt, left=8pt, right=8pt, top=6pt, bottom=6pt,
  title={\textbf{TravelShare — Réseau social de voyages}},
  fonttitle=\bfseries\small\color{white},
  colbacktitle=teal]

\begin{multicols}{2}
\small
\textbf{Compte \& Session}
\begin{itemize}[leftmargin=*,itemsep=0pt,topsep=2pt]
  \item Inscription / connexion avec session persistante
  \item Mode anonyme (consultation des photos publiques)
  \item Profil public cliquable (style Instagram)
  \item Photo de profil + bio utilisateur
\end{itemize}

\textbf{Publication}
\begin{itemize}[leftmargin=*,itemsep=0pt,topsep=2pt]
  \item Photo depuis galerie ou caméra
  \item Géocodage automatique via Nominatim
  \item Note vocale AAC (enregistrement + lecture)
  \item Tags auto + IA ML Kit Image Labeling (seuil 70\,\%)
  \item Localisation approximative (±11 km)
  \item Visibilité PUBLIC ou GROUP (groupe privé)
\end{itemize}

\textbf{Exploration \& Recherche}
\begin{itemize}[leftmargin=*,itemsep=0pt,topsep=2pt]
  \item Recherche texte temps réel + recherche vocale
  \item 8 filtres : catégorie, auteur, période, GPS, likes, aléatoire
  \item Chips en 2 rangées fixes (pas de défilement)
  \item Pull-to-refresh + pagination infinite scroll
  \item Shake pour flux aléatoire (accéléromètre)
\end{itemize}

\columnbreak

\textbf{Fiche photo}
\begin{itemize}[leftmargin=*,itemsep=0pt,topsep=2pt]
  \item Likes persistés, commentaires, signalement
  \item Temps de marche via OSRM (profil piéton)
  \item Photos similaires par tags
  \item Partage externe (WhatsApp, SMS…)
  \item Partage dans un groupe (carte cliquable dans le chat)
  \item Modifier / supprimer sa photo ou ses commentaires
\end{itemize}

\textbf{Groupes \& Chat}
\begin{itemize}[leftmargin=*,itemsep=0pt,topsep=2pt]
  \item Groupes privés : création, demande PENDING→MEMBER
  \item Chat temps réel (LiveData Room) — 3 types de messages
  \item Badges non lus dynamiques par groupe
\end{itemize}

\textbf{Notifications}
\begin{itemize}[leftmargin=*,itemsep=0pt,topsep=2pt]
  \item Notifications système Android
  \item Notifications in-app + badge cloche
  \item Préférences configurables (AUTEUR / LIEU / TAG / GROUPE)
\end{itemize}

\textbf{Carte}
\begin{itemize}[leftmargin=*,itemsep=0pt,topsep=2pt]
  \item Carte OSMDroid avec tous les points géolocalisés
\end{itemize}
\end{multicols}
\end{tcolorbox}

\vspace{8pt}

% ── TravelPath ───────────────────────────────────────────────────
\begin{tcolorbox}[enhanced, colback=lightbg, colframe=navy, arc=4pt,
  boxrule=0.6pt, left=8pt, right=8pt, top=6pt, bottom=6pt,
  title={\textbf{TravelPath — Générateur d'itinéraires}},
  fonttitle=\bfseries\small\color{white},
  colbacktitle=navy]

\begin{multicols}{2}
\small
\textbf{Formulaire de génération}
\begin{itemize}[leftmargin=*,itemsep=0pt,topsep=2pt]
  \item Ville, 5 activités, budget, durée, effort, météo
  \item Lieux obligatoires à intégrer dans le parcours
  \item Génération de 3 parcours : Éco / Équilibré / Confort
\end{itemize}

\textbf{Vrais lieux OSM}
\begin{itemize}[leftmargin=*,itemsep=0pt,topsep=2pt]
  \item Overpass API — lieux nommés uniquement
  \item 5 requêtes parallèles (1 par activité)
  \item Fallback Nominatim si Overpass indisponible
  \item Déduplication des lieux identiques
  \item Adresse réelle de chaque étape
\end{itemize}

\textbf{Cartographie}
\begin{itemize}[leftmargin=*,itemsep=0pt,topsep=2pt]
  \item Carte OSMDroid avec marqueurs par étape
  \item Tracé routier réel OSRM + temps de trajet
  \item Géocodage Nominatim + validation Haversine
  \item Zoom / déplacement libres sur la carte
\end{itemize}

\columnbreak

\textbf{Métriques dynamiques}
\begin{itemize}[leftmargin=*,itemsep=0pt,topsep=2pt]
  \item Budget, durée et effort pris en compte
  \item Coûts proportionnels (25\,\% / 60\,\% / 90\,\%)
  \item Météo temps réel Open-Meteo
  \item Horaires d'ouverture typiques par type
\end{itemize}

\textbf{Gestion des parcours}
\begin{itemize}[leftmargin=*,itemsep=0pt,topsep=2pt]
  \item Onglets Mes parcours / Sauvegardés
  \item Suppression individuelle ou globale
  \item Régénération avec paramètres pré-remplis
  \item Export PDF natif (Canvas A4)
  \item Partage texte du parcours
  \item Mode hors-ligne (plans sauvegardés Room)
\end{itemize}

\textbf{Détail des étapes}
\begin{itemize}[leftmargin=*,itemsep=0pt,topsep=2pt]
  \item Icône, créneau horaire, durée, coût
  \item Bouton vidéo YouTube par étape
\end{itemize}
\end{multicols}
\end{tcolorbox}

\vspace{8pt}

% ── Passerelle ───────────────────────────────────────────────────
\begin{tcolorbox}[enhanced, colback=okbg, colframe=greentxt, arc=4pt,
  boxrule=0.6pt, left=8pt, right=8pt, top=6pt, bottom=6pt,
  title={\textbf{Passerelle — Intégration bidirectionnelle}},
  fonttitle=\bfseries\small\color{white},
  colbacktitle=greentxt]
\small
\begin{multicols}{2}
\begin{itemize}[leftmargin=*,itemsep=0pt,topsep=2pt]
  \item TravelShare → TravelPath : ville transmise automatiquement
  \item TravelPath → TravelShare : bouton "Voir des photos" par étape
\end{itemize}
\columnbreak
\begin{itemize}[leftmargin=*,itemsep=0pt,topsep=2pt]
  \item Navigation unifiée (BottomNavigationView + back stack)
  \item Session partagée entre les deux modules
\end{itemize}
\end{multicols}
\end{tcolorbox}

% ==============================================================
\chapter{Présentation générale de l'application}
% ==============================================================

\section{Architecture globale : le pattern MVVM}

Toute l'application est construite sur le pattern \textbf{MVVM
(Model-View-ViewModel)}, le standard recommandé par Google pour Android.

L'idée est simple : la couche d'affichage (les Fragments et Activities) ne
sait pas d'où viennent les données. Elle se contente d'\emph{observer} des
\ic{LiveData} que lui fournissent les ViewModels. Quand les données changent
dans la base, l'écran se met à jour automatiquement, sans qu'il faille écrire
de code de synchronisation.

\begin{tcolorbox}[mybox,colback=codebg,colframe=bord,
  left=10pt,right=10pt,top=8pt,bottom=8pt]
{\ttfamily\footnotesize\color{navy!80}
\begin{verbatim}
 ┌─────────────── COUCHE VUE (Fragments / Activities) ──────────────┐
 │  MainActivity · ExplorerFragment · PublishFragment               │
 │  GroupsFragment · GroupChatFragment · MapFragment                │
 │  TravelPathFragment · PlanDetailFragment · ...                   │
 └────────────────────────┬─────────────────────────────────────────┘
                          │  observe(LiveData)  /  appel méthode VM
                          ▼
 ┌─────────────── COUCHE VIEWMODEL ─────────────────────────────────┐
 │  SharedViewModel              TravelPathViewModel                │
 │  (photos, groupes, notifs)    (parcours, étapes, APIs)           │
 └────────────────────────┬─────────────────────────────────────────┘
                          │  appels DAO sur thread background
                          ▼
 ┌─────────────── COUCHE MODÈLE (Room v19) ──────────────────────────┐
 │  AppDatabase · 12 entités · pool de 4 threads                    │
 │  PhotoDao · GroupDao · TravelPlanDao · PlanStepDao · ...         │
 └────────────────────────┬─────────────────────────────────────────┘
                          │  requêtes HTTP (thread background)
                          ▼
 ┌─────────────── APIs EXTERNES (gratuites) ─────────────────────────┐
 │  Nominatim (géocodage)  ·  OSRM (routage)  ·  Open-Meteo (météo) │
 └───────────────────────────────────────────────────────────────────┘
\end{verbatim}}
\end{tcolorbox}


\section{La base de données : Room v19}

Room est la librairie officielle Android pour gérer une base SQLite de façon
sûre et structurée. Notre base s'appelle \ic{tp3\_database} et contient
\textbf{12 entités} (tables). Chaque entité est une classe Java annotée avec
\ic{@Entity}, et chaque DAO (Data Access Object) est une interface annotée avec
\ic{@Dao} qui définit les requêtes SQL.

\begin{tcolorbox}[enhanced,colback=lightbg,colframe=bord,arc=5pt,boxrule=0.5pt,
  left=8pt,right=8pt,top=8pt,bottom=8pt]
\textbf{Exemple concret :} quand l'utilisateur tape dans la barre de recherche,
un \ic{TextWatcher} appelle \ic{viewModel.searchPhotos("paris")}. Le ViewModel
appelle \ic{photoDao.searchPhotos("paris")} qui exécute une requête SQL
\ic{LIKE '\%paris\%'} et retourne un \ic{LiveData<List<Photo>>}. Room notifie
automatiquement l'UI, qui met à jour la liste de photos.
\end{tcolorbox}

Toutes les opérations de base de données s'exécutent sur un pool de \textbf{4 threads background} (\ic{ExecutorService}), jamais sur le thread principal.

\section{Navigation et structure de l'application}

La navigation principale utilise une \ic{BottomNavigationView} avec 5 onglets :
Explorateur, Carte, Publier, Groupes, Profil. Chaque onglet charge un Fragment
dans un conteneur central. Les sous-écrans (détail d'un parcours, chat de groupe...)
s'empilent par-dessus via \ic{addToBackStack()}, permettant de revenir en arrière.

\screenshot{screen_main.png}

\section{Choix techniques justifiés}

\begin{tabularx}{\textwidth}{>{\bfseries}l X}
  \toprule
  Choix & Pourquoi ce choix \\
  \midrule
  Room (SQLite local)
    & Les données de l'utilisateur (photos, groupes, parcours) sont stockées
      localement. Room offre la sécurité de type (requêtes vérifiées à la compilation),
      les LiveData automatiques, et fonctionne sans connexion réseau. \\
  OSMDroid (cartographie)
    & Alternative open-source à Google Maps, sans clé API ni facturation.
      Supporte les marqueurs, les polylines et le zoom automatique. \\
  Nominatim + OSRM + Open-Meteo + Overpass
    & Quatre services entièrement gratuits et sans clé. Nominatim géocode les
      adresses et fait le reverse geocoding, OSRM calcule les itinéraires routiers,
      Open-Meteo fournit la météo, Overpass retourne les vrais POIs OSM
      (restaurants, musées, parcs…) avec leur nom et coordonnées exactes. \\
  ActivityResultLauncher
    & API moderne (remplace \ic{onActivityResult} deprecated depuis API 30)
      pour la caméra, la galerie et les permissions runtime. \\
  FileProvider
    & Obligatoire depuis Android 7 pour partager des fichiers entre applications.
      Utilisé pour la caméra, les enregistrements audio et l'export PDF. \\
  ML Kit Image Labeling
    & SDK Google on-device (gratuit, sans clé API). Analyse les photos sélectionnées
      et retourne des labels visuels (ex. "Mountain", "Sky") avec un score de
      confiance. Seuil fixé à 70\,\% pour éviter les faux positifs. \\
  SwipeRefreshLayout
    & Composant Jetpack ajouté dans l'explorateur. Détecte le geste "tirer vers le bas"
      et déclenche un rechargement de la première page de photos publiques. \\
  \bottomrule
\end{tabularx}


% ==============================================================
\chapter{Module TravelShare — Partage de photos}
% ==============================================================

\section{Présentation et modes d'utilisation}

TravelShare est le module de réseau social de l'application. Il est organisé
autour de deux niveaux d'accès distincts selon que l'utilisateur dispose ou
non d'un compte.

\textbf{Mode anonyme} : l'utilisateur peut parcourir et filtrer toutes les
photos publiques, consulter le détail d'une photo (lieu, date, commentaires,
coordonnées GPS). Il peut liker et signaler un contenu, mais toute action
d'écriture est bloquée avec un message d'invite à la connexion.

\textbf{Mode connecté} : toutes les fonctionnalités sont disponibles.
L'utilisateur publie des photos, ajoute une note vocale, rejoint des groupes
privés, discute en chat, reçoit des notifications personnalisées et configure
ses préférences d'alerte.

La session est gérée par \ic{SessionManager}, qui lit et écrit dans
\ic{SharedPreferences}. Elle stocke l'\ic{userId} et le \ic{username}
de l'utilisateur connecté. Chaque fragment interroge \ic{session.isLoggedIn()}
pour décider des actions autorisées.

\section{Authentification (LoginActivity et InscriptionFragment)}

\ic{LoginActivity} est l'activité de lancement de l'application. Elle affiche
directement l'écran de connexion sous le titre \textbf{Traveling}, avec trois
actions possibles : se connecter, créer un compte, ou continuer en mode anonyme.
Aucune page d'accueil intermédiaire n'existe — l'utilisateur arrive immédiatement
sur cet écran au démarrage.

La connexion vérifie les identifiants avec \ic{userDao.getUserByLogin()}. Une
fois connecté, \ic{SessionManager.createLoginSession()} enregistre l'\ic{userId}
et le \ic{username} en \ic{SharedPreferences}. Si l'utilisateur était déjà
connecté lors du lancement, \ic{LoginActivity} le redirige directement vers
\ic{MainActivity} sans afficher le formulaire.

Depuis \ic{MainActivity}, si l'Intent contient \ic{OPEN\_FRAGMENT = "INSCRIPTION"},
la \ic{BottomNavigationView} est masquée et \ic{InscriptionFragment} est chargé
directement. Quand l'inscription réussit, le callback \ic{onInscriptionReussie()}
de l'interface \ic{OnInscriptionListener} est déclenché, la session est créée
et l'utilisateur est redirigé vers l'explorateur.

\section{L'explorateur de photos (ExplorerFragment)}

L'explorateur est le fragment principal chargé au démarrage. Il affiche un
\ic{RecyclerView} en grille à 2 colonnes via \ic{GridLayoutManager}, alimenté
par un \ic{PhotoAdapter} qui charge les images avec Glide. Il propose huit
façons distinctes de filtrer les photos publiques, toutes gérées par la même
méthode \ic{observeSource()} qui garantit qu'une seule source LiveData est
active à la fois.

\begin{lstlisting}
private LiveData<List<Photo>> currentSource;

private void observeSource(LiveData<List<Photo>> newSource) {
    if (currentSource != null)
        currentSource.removeObservers(getViewLifecycleOwner());
    currentSource = newSource;
    currentSource.observe(getViewLifecycleOwner(), adapter::setPhotos);
}
\end{lstlisting}

\begin{tabularx}{\textwidth}{>{\bfseries}l X}
  \toprule
  Filtre & Implémentation \\
  \midrule
  Recherche texte
    & \ic{TextWatcher} sur un \ic{EditText}. À chaque modification de texte,
      appelle \ic{viewModel.searchPhotos(q)} qui exécute une requête SQL
      \ic{LIKE '\%q\%'} sur le titre, le lieu et les tags.
      Si le champ est vide, revient sur \ic{getPublicPhotos()}. \\
  Recherche vocale
    & Bouton micro qui lance un \ic{RecognizerIntent.ACTION\_RECOGNIZE\_SPEECH}.
      Le résultat texte est récupéré dans \ic{onActivityResult(REQUEST\_VOICE)}
      et injecté directement dans l'\ic{EditText} de recherche, ce qui
      déclenche automatiquement le \ic{TextWatcher}. \\
  Filtre par auteur
    & Second \ic{EditText} avec son propre \ic{TextWatcher}. Appelle
      \ic{getPublicPhotosByAuthor(author)}. Si vide, revient sur toutes
      les photos publiques. \\
  Filtre par période
    & Deux \ic{EditText} (date début, date fin). Un \ic{OnFocusChangeListener}
      commun se déclenche quand l'utilisateur quitte l'un des deux champs.
      Si les deux sont renseignés, appelle \ic{getPhotosByDateRange(start, end)}. \\
  Chips de catégorie
    & Huit \ic{TextView} stylés en pill, disposés en \textbf{deux rangées fixes
      de quatre} (pas de défilement horizontal). Chaque chip occupe la même
      largeur via \ic{layout\_weight="1"}. Le chip actif passe en fond solide
      et texte blanc via \ic{setActiveChip(chipId)}.
      Appelle \ic{getPhotosByCategory(category)} ou \ic{getPublicPhotos()} pour "Tous". \\
  Autour de moi (GPS)
    & Vérifie la permission \ic{ACCESS\_FINE\_LOCATION} via
      \ic{ActivityResultLauncher}. Récupère la position avec
      \ic{LocationManager.getLastKnownLocation()} (GPS puis réseau en fallback).
      Appelle \ic{getPhotosByLocation(lat, lng, delta)} où \ic{delta = 50/111}
      correspond à un rayon de 50 km. \\
  ❤️ Mes Likes
    & Affiche toutes les photos que l'utilisateur a likées, dans n'importe
      quelle session. Les IDs likés sont persistés en \ic{SharedPreferences}
      (clé \ic{liked\_ids}). Le chip appelle \ic{viewModel.getPhotosByIds(ids)}
      qui exécute un \ic{SELECT ... WHERE id IN (:ids)} en Room. Si la liste
      est vide, un toast invite à liker des photos. \\
  Flux aléatoire
    & Bouton dédié et détection de shake (accéléromètre). Appelle
      \ic{getRandomPhotos(10)} qui exécute un \ic{SELECT ... ORDER BY RANDOM() LIMIT 10}. \\
  \bottomrule
\end{tabularx}

\subsection{Pull-to-refresh et pagination}

L'explorateur intègre un \ic{SwipeRefreshLayout} encapsulant le \ic{RecyclerView}.
Tirer vers le bas recharge la première page de photos publiques et réinitialise
le curseur de pagination.

La liste principale (chip "Tous") fonctionne en mode paginé : 20 photos sont
chargées initialement via \ic{photoDao.getPublicPhotosPage(offset, limit)}.
Un \ic{RecyclerView.OnScrollListener} surveille la position du dernier élément visible.
Quand il atteint les 4 derniers items, la page suivante est chargée en background
et ajoutée à l'adapter via \ic{appendPhotos()} sans \ic{notifyDataSetChanged()}.

\begin{lstlisting}
recyclerView.addOnScrollListener(new RecyclerView.OnScrollListener() {
    @Override
    public void onScrolled(@NonNull RecyclerView rv, int dx, int dy) {
        if (!isPaginationMode || dy <= 0 || isLoadingMore) return;
        int lastVisible = layoutManager.findLastVisibleItemPosition();
        if (lastVisible >= adapter.getItemCount() - 4)
            loadPage(currentOffset, false); // false = append
    }
});
\end{lstlisting}

Les vues filtrées (recherche, catégorie, GPS…) utilisent toujours \ic{LiveData}
directement — la pagination n'est active que sur le flux public non filtré.

\screenshot{screen_explorer.png}

\section{Capteur accéléromètre : détection du shake}

\ic{ExplorerFragment} implémente \ic{SensorEventListener}. En secouant le
téléphone, un flux de 10 photos aléatoires s'affiche. L'enregistrement du
capteur se fait dans \ic{onResume()} et le désenregistrement dans \ic{onPause()}
pour économiser la batterie.

Le shake est détecté en calculant la variation d'accélération sur les trois axes
entre deux mesures consécutives. Si la variation dépasse \ic{SHAKE\_THRESHOLD}
sur au moins deux axes simultanément, et que le cooldown de 1500 ms est écoulé,
le shake est validé.

\begin{lstlisting}
private static final float SHAKE_THRESHOLD   = 12f;
private static final long  SHAKE_COOLDOWN_MS = 1500;

@Override
public void onSensorChanged(SensorEvent event) {
    float deltaX = Math.abs(event.values[0] - lastX);
    float deltaY = Math.abs(event.values[1] - lastY);
    float deltaZ = Math.abs(event.values[2] - lastZ);
    lastX = event.values[0]; lastY = event.values[1]; lastZ = event.values[2];

    boolean shook = (deltaX > SHAKE_THRESHOLD && deltaY > SHAKE_THRESHOLD)
                 || (deltaX > SHAKE_THRESHOLD && deltaZ > SHAKE_THRESHOLD)
                 || (deltaY > SHAKE_THRESHOLD && deltaZ > SHAKE_THRESHOLD);

    if (shook) {
        long now = System.currentTimeMillis();
        if (now - lastShakeTime > SHAKE_COOLDOWN_MS) {
            lastShakeTime = now;
            setActiveChip(-1); // aucun chip selectionne
            observeSource(viewModel.getRandomPhotos(10));
        }
    }
}
\end{lstlisting}

\section{Publication de photos (PublishFragment)}

\ic{PublishFragment} permet à un utilisateur connecté de publier une photo avec
l'ensemble de ses métadonnées. Le formulaire comprend :

\begin{itemize}
  \item \textbf{Sélection de l'image} : un clic sur la zone d'aperçu ouvre un
    \ic{AlertDialog} proposant "Galerie" ou "Appareil photo". La galerie utilise
    \ic{galleryLauncher.launch("image/*")} (\ic{ActivityResultContracts.GetContent}).
    La caméra crée d'abord un fichier dans \ic{externalCacheDir/images/} via
    \ic{FileProvider}, puis lance \ic{cameraLauncher.launch(uri)}
    (\ic{ActivityResultContracts.TakePicture}). L'URI résultante est stockée
    dans \ic{selectedPhotoUri} et affichée dans \ic{ImageView}.

  \item \textbf{Titre et lieu} : deux \ic{EditText} obligatoires. Le lieu est
    géocodé en arrière-plan via Nominatim au moment de la publication, ce qui
    permet d'afficher la photo sur la carte avec les bonnes coordonnées GPS.

  \item \textbf{Catégorie} : un \ic{Spinner} avec quatre valeurs (Nature,
    Urbain, Culture, Magasin), alimenté par un \ic{ArrayAdapter}.

  \item \textbf{Tags automatiques + IA (ML Kit)} : dès qu'une photo est
    sélectionnée (galerie ou caméra), \ic{runMlKit(uri)} analyse l'image en
    arrière-plan avec \ic{ImageLabeling.getClient()} (seuil 70\,\%). Quand
    l'utilisateur clique ensuite sur "Tags auto", les hashtags mots-clés
    (lieu, catégorie, titre) sont fusionnés avec les labels ML Kit
    (\ic{\#Mountain}, \ic{\#Sky}…). Un toast "✨ X tags IA ajoutés"
    confirme l'enrichissement IA.

  \item \textbf{Note vocale} : décrit en détail ci-dessous.

  \item \textbf{Localisation approximative} : une \ic{CheckBox}
    "📍 Lieu approximatif" permet de masquer la position exacte. Si cochée,
    les coordonnées issues de Nominatim sont arrondies à 1 décimale
    (\texttt{Math.round(lat * 10.0) / 10.0}), ce qui correspond à une
    précision de ±11 km. La fiche photo affichera alors "≈ XX.X° N"
    au lieu de la latitude exacte.

  \item \textbf{Visibilité} : un \ic{Switch} bascule entre PUBLIC et GROUP.
    En mode GROUP, un \ic{Spinner} chargé depuis \ic{viewModel.getGroupsForUser()}
    permet de sélectionner le groupe destinataire. L'utilisateur peut aussi
    créer un nouveau groupe directement depuis ce formulaire.
\end{itemize}

Au moment de la publication, un \ic{Photo} est créé avec toutes les
métadonnées (y compris \ic{voiceUri} si une note vocale a été enregistrée),
inséré en base via \ic{viewModel.insert(photo)}, puis
\ic{triggerNotificationsForPublish()} est appelé pour envoyer les
notifications aux utilisateurs abonnés correspondants.

\screenshot{screen_publish.png}

\subsection{La note vocale (MediaRecorder)}

Le bouton de note vocale fonctionne selon une machine d'état à trois états.
Le même bouton change de rôle à chaque appui.

\begin{lstlisting}
private enum VoiceState { IDLE, RECORDING, RECORDED }
// IDLE    -> appui -> demande permission RECORD_AUDIO -> RECORDING
// RECORDING -> appui -> stopVoiceRecording() -> RECORDED
// RECORDED  -> appui -> playVoiceNote() via MediaPlayer -> reste RECORDED
\end{lstlisting}

L'enregistrement utilise \ic{MediaRecorder} en format MPEG\_4 / codec AAC,
44100 Hz, 128 kbps, mono. Le fichier est sauvegardé dans
\ic{externalCacheDir/voice/voice\_TIMESTAMP.m4a}. En fin d'enregistrement,
\ic{voiceNotePath} est associé à la photo lors de la publication.
Dans \ic{onDestroyView()}, le \ic{MediaRecorder} et le \ic{MediaPlayer} sont
libérés pour éviter les fuites de ressources.

\begin{lstlisting}
private void startVoiceRecording(Button btn) {
    mediaRecorder = new MediaRecorder();
    mediaRecorder.setAudioSource(MediaRecorder.AudioSource.MIC);
    mediaRecorder.setOutputFormat(MediaRecorder.OutputFormat.MPEG_4);
    mediaRecorder.setAudioEncoder(MediaRecorder.AudioEncoder.AAC);
    mediaRecorder.setAudioSamplingRate(44100);
    mediaRecorder.setAudioEncodingBitRate(128000);
    mediaRecorder.setAudioChannels(1);
    mediaRecorder.setOutputFile(voiceNotePath);
    mediaRecorder.prepare();
    mediaRecorder.start();
    voiceState = VoiceState.RECORDING;
}
\end{lstlisting}

La permission \ic{RECORD\_AUDIO} est demandée via \ic{ActivityResultLauncher}
au premier appui. Si l'utilisateur accepte, \ic{startVoiceRecording()} est
appelé immédiatement depuis le callback de la permission.

\section{Consultation d'une photo (PhotoDetailActivity)}

En cliquant sur une photo dans la grille, \ic{PhotoDetailActivity} s'ouvre.
Elle affiche l'image en grand (chargée via Glide), le titre, le lieu, l'auteur,
la date, la catégorie, les tags et les coordonnées GPS.

Depuis cet écran, l'utilisateur peut :
\begin{itemize}
  \item \textbf{Liker / Unliker} : incrémente ou décrémente le compteur de likes
    via \ic{viewModel.updateLikes(photoId, newCount)}. Le bouton se met à jour
    immédiatement sans rechargement.
  \item \textbf{Commenter} : un \ic{EditText} + bouton "Envoyer" insère un
    \ic{Comment} en base via \ic{viewModel.insertComment(comment)}. La liste
    des commentaires, chargée depuis \ic{getCommentsForPhoto(photoId)} en
    LiveData, se met à jour automatiquement.
  \item \textbf{Signaler} : insère un \ic{Report} en base via
    \ic{viewModel.insertReport(report)}. Le bouton est désactivé après
    l'envoi pour éviter les doublons.
  \item \textbf{Naviguer vers le lieu} : un bouton lance un \ic{Intent} avec
    l'URI \ic{geo:LAT,LNG?q=NOM\_DU\_LIEU}, qui ouvre Google Maps (ou toute
    application de navigation installée).
  \item \textbf{Partager} : le bouton "↗ Partager" construit un résumé texte
    (titre, lieu, catégorie, tags) et le propose via \ic{Intent.ACTION\_SEND}
    à n'importe quelle application (WhatsApp, SMS, e-mail…).
  \item \textbf{Partager dans un groupe} : le bouton "👥 Partager dans un groupe"
    charge en arrière-plan la liste des groupes dont l'utilisateur est créateur
    ou membre admis (\ic{groupDao.getGroupsForUserSync(userId)}). Un
    \ic{AlertDialog} permet de choisir le groupe cible. Un \ic{GroupMessage}
    est inséré avec le champ \ic{photoId} renseigné, ce qui permet au chat
    du groupe d'afficher une carte "🔗 Photo partagée" cliquable distincte
    des messages texte et des posts publiés directement dans le groupe.
  \item \textbf{Modifier le titre} (si auteur) : si l'utilisateur connecté est
    l'auteur de la photo, un bloc "Actions propriétaire" apparaît. Le bouton
    "✏ Modifier le titre" ouvre un \ic{AlertDialog} avec un \ic{EditText}
    pré-rempli. La validation appelle \ic{viewModel.updatePhotoTitle(photoId, newTitle)}.
  \item \textbf{Supprimer sa photo} (si auteur) : le bouton "🗑 Supprimer" ouvre
    une boîte de confirmation. Si confirmé, \ic{viewModel.deletePhoto(photoId)} est
    appelé et l'activité se ferme.
  \item \textbf{Supprimer ses commentaires} : le \ic{CommentAdapter} reçoit
    l'\ic{userId} courant. Si \ic{comment.userId == currentUserId}, un bouton "×"
    terracotta apparaît sur le commentaire ; son clic appelle
    \ic{viewModel.deleteComment(id)}.
\end{itemize}

\subsection{Note vocale attachée à la photo}

Si la photo possède un \ic{voiceUri} (enregistrement M4A sauvegardé lors de la
publication), un bloc "🎙 Note vocale" apparaît avec un bouton ▶/⏸. La lecture
utilise \ic{MediaPlayer} avec gestion d'état (\ic{isPlaying}) et libération des
ressources dans \ic{onDestroy()} pour éviter les fuites mémoire.

\subsection{Temps de marche vers la photo}

Le bloc ACCÈS affiche dynamiquement le temps de marche depuis la position actuelle
de l'utilisateur. Au chargement, la permission \ic{ACCESS\_FINE\_LOCATION} est
vérifiée via \ic{ActivityResultLauncher}. Si accordée, \ic{LocationManager.getLastKnownLocation()}
récupère la position GPS, puis une requête HTTP appelle OSRM avec le profil
\ic{foot} (marche à pied) :

\begin{lstlisting}
// Format : .../route/v1/foot/{userLng},{userLat};{photoLng},{photoLat}
// Reponse : duration (secondes) + distance (metres)
// Affichage : "🚶 12 min (950 m)" ou "🚶 1h05 (5.2 km)"
\end{lstlisting}

En cas d'échec réseau ou de coordonnées manquantes, le libellé revient à "Maps →".

\screenshot{screen_photo_detail.png}

\subsection{Photos similaires par tags}

La section "PHOTOS SIMILAIRES" recherche d'abord les photos partageant le
\textbf{premier tag} de la photo courante via \ic{viewModel.getPhotosByTag(tag)}.
La méthode \ic{extractFirstTag()} parse la chaîne de tags (\ic{\#Paris \#Voyage …})
et extrait le premier mot de plus de 2 caractères. Si aucun résultat par tag,
la recherche retombe sur \ic{getPhotosByCategory(category)}.

\section{Les groupes (GroupsFragment)}

\ic{GroupsFragment} est organisé en deux onglets gérés manuellement :

\textbf{Onglet "Mes groupes"} : affiche uniquement les groupes dont l'utilisateur
est créateur ou membre accepté, via \ic{viewModel.getGroupsForUser(userId)}.
Si l'utilisateur n'est pas connecté, la liste est vide (pas d'appel à
\ic{getAllGroups()}).

\textbf{Onglet "Découvrir"} : affiche tous les groupes via
\ic{viewModel.getAllGroups()}, avec une barre de recherche qui filtre via
\ic{viewModel.searchGroups(query)}. Pour chaque groupe affiché, le statut
de l'utilisateur courant (null, PENDING, MEMBER) est résolu en parallèle
par un appel \ic{viewModel.getMembership()} sur le
\ic{databaseWriteExecutor}, puis mis à jour sur l'UI via \ic{runOnUiThread()}.

Le \ic{GroupAdapter} gère trois états d'affichage pour le bouton d'action :
\begin{itemize}
  \item \textbf{null} (non membre) : bouton "Demander →" — appelle
    \ic{requestJoinGroup()} qui insère un \ic{GroupMember} avec statut
    \ic{PENDING}, et envoie une \ic{AppNotification} au créateur.
  \item \textbf{PENDING} : bouton "⏳ En attente" non cliquable.
  \item \textbf{MEMBER} : bouton "Membre ✓" — clic propose de quitter le groupe
    via \ic{rejectOrLeaveGroup()}.
\end{itemize}

Pour créer un groupe, un \ic{AlertDialog} demande nom + description, puis
appelle \ic{viewModel.insertGroup(group)} et bascule automatiquement sur
l'onglet "Mes groupes".

\screenshot{screen_groups.png}

\subsection{Chat de groupe (GroupChatFragment)}

\ic{GroupChatFragment} est ouvert depuis l'onglet "Mes groupes" en cliquant
sur "Chat". Il reçoit le \ic{groupId}, le \ic{groupName} et le
\ic{creatorId} en arguments.

Le chat fusionne trois types d'éléments dans un seul flux trié par date,
via un \ic{MediatorLiveData} combinant messages et posts :

\begin{tabularx}{\textwidth}{>{\bfseries}l X}
  \toprule
  Type & Affichage \\
  \midrule
  Message texte (\ic{TYPE\_MESSAGE})
    & Bulle teal (messages envoyés par l'utilisateur) ou sable (autres).
      Alignée à droite ou à gauche selon l'auteur. \\
  Post publié dans le groupe (\ic{TYPE\_POST})
    & Carte blanche avec barre teal, miniature de la photo si disponible,
      titre et lieu. Cliquable vers \ic{PhotoDetailActivity}. \\
  Photo partagée (\ic{TYPE\_SHARED})
    & Carte distincte avec bandeau \textbf{terracotta} "🔗 Photo partagée",
      miniature chargée depuis la base via \ic{photoId}, mention "Voir la photo →".
      Détectée via \ic{GroupMessage.photoId > 0}. Cliquable vers \ic{PhotoDetailActivity}. \\
  \bottomrule
\end{tabularx}

L'envoi d'un message texte crée un \ic{GroupMessage} (champ \ic{photoId = 0})
et appelle \ic{viewModel.sendGroupMessage(msg)}.

\screenshot{screen_chat.png}

\subsection{Badges de messages non lus}

Dans l'onglet "Mes groupes", chaque item affiche un badge rouge circulaire
superposé à l'avatar du groupe. Ce badge est alimenté par une requête
\ic{COUNT} sur la table \ic{AppNotification} :

\begin{lstlisting}
viewModel.getUnreadCountForGroup(session.getUserId(), g.id)
    .observe((LifecycleOwner) fragment, count -> {
        if (count != null && count > 0) {
            h.tvUnreadBadge.setVisibility(View.VISIBLE);
            h.tvUnreadBadge.setText(count > 9 ? "9+" : count.toString());
        } else {
            h.tvUnreadBadge.setVisibility(View.GONE);
        }
    });
\end{lstlisting}

Quand l'utilisateur ouvre le chat, \ic{markGroupMessagesRead(userId, groupId)}
met à jour \ic{isRead = 1} pour toutes les notifications de ce groupe,
et le badge disparaît immédiatement via la LiveData.

\section{Carte interactive (MapFragment)}

\ic{MapFragment} utilise OSMDroid pour afficher sur une carte tous les points
de photos géolocalisées. Il observe \ic{viewModel.getPublicPhotos()} et
place un marqueur pour chaque photo dont lat et lng sont non nuls.
Un clic sur un marqueur affiche une info-bulle avec le titre et le lieu,
et un second clic ouvre le détail de la photo.

\screenshot{screen_map.png}

\section{Système de notifications}

Le système de notifications fonctionne sur deux niveaux indépendants et
complémentaires.

\textbf{Notifications système Android} : après chaque publication,
\ic{triggerNotificationsForPublish()} charge les préférences de l'utilisateur
(\ic{prefDao.getPreferencesForUserSync(userId)}) et vérifie si chaque
préférence correspond à la publication. Si oui, \ic{NotificationUtil.showNotification()}
crée une notification dans le canal \ic{TravelShare} qui apparaît dans la
barre de statut Android.

\textbf{Notifications in-app} : stockées dans la table \ic{AppNotification}.
Chaque type d'événement (JOIN\_REQUEST, GROUP\_MESSAGE, etc.) insère une entrée
via \ic{viewModel.insertAppNotification(notif)}. \ic{NotificationsFragment}
observe \ic{getAppNotificationsForUser(userId)} et les affiche dans une liste.
L'utilisateur peut marquer une notification comme lue individuellement, ou
tout effacer. Un badge rouge sur la cloche de \ic{MainActivity} compte les
non-lues via \ic{getUnreadNotificationCount(userId)}.

\textbf{Préférences configurables} : dans \ic{NotificationPreferencesFragment},
l'utilisateur définit sur quels événements il veut être alerté via quatre types :
\ic{AUTEUR} (auteur spécifique), \ic{LIEU} (lieu contenant la valeur),
\ic{TAG} (tag correspondant) et \ic{GROUPE} (catégorie correspondante).


\section{Photo de profil et bio (ProfileFragment)}

\ic{ProfileFragment} a été enrichi d'un avatar photo et d'une biographie
utilisateur, stockés dans la table \ic{User} (champs \ic{avatarUri} et \ic{bio}).

L'avatar est un \ic{FrameLayout} superposant un \ic{TextView} (initiale de
l'utilisateur, fallback) et un \ic{ImageView} (photo chargée avec Glide).
Un clic sur ce bloc ouvre le sélecteur d'images système via
\ic{ActivityResultContracts.GetContent("image/*")}. L'URI retournée est persistée
dans la base via \ic{userDao.updateUserProfile(userId, avatarUri, bio)}.

La biographie est un \ic{EditText} pré-rempli depuis la base au chargement du
fragment. Un bouton "OK" déclenche la sauvegarde. Le chargement initial se fait
en background (thread pool) avec retour sur l'UI via \ic{runOnUiThread()}.

\section{Profil public d'un utilisateur (UserProfileActivity)}

En cliquant sur le nom d'un auteur dans l'explorateur ou sur la fiche d'une
photo, l'utilisateur accède au profil public de cet auteur via
\ic{UserProfileActivity}. Cette vue, inspirée du modèle Instagram, présente :

\begin{itemize}
  \item Un avatar (initiale ou photo de profil chargée via Glide), le nom
    d'utilisateur, et le compteur de publications.
  \item La biographie si elle a été renseignée.
  \item Une grille à 3 colonnes (\ic{GridLayoutManager}) de toutes les photos
    publiques de cet auteur, chargées via \ic{getPublicPhotosByAuthor(username)}
    (LiveData Room). Cliquer sur une photo ouvre \ic{PhotoDetailActivity}.
\end{itemize}

Le déclenchement depuis \ic{PhotoAdapter} s'effectue par un \ic{setOnClickListener}
sur le \ic{TextView} \ic{@auteur} de chaque carte. Depuis \ic{PhotoDetailActivity},
c'est la ligne "date · par @auteur" qui est cliquable. Dans les deux cas,
un \ic{Intent} est lancé avec \ic{EXTRA\_USERNAME} pour ouvrir le bon profil.

\screenshot{screen_user_profile.png}

% ==============================================================
\chapter{Module TravelPath — Générateur d'itinéraires}
% ==============================================================

\section{Présentation}

TravelPath génère des itinéraires de visite personnalisés à partir des
préférences de l'utilisateur. Il produit jusqu'à 3 parcours thématiques
(économique, équilibré, confort), les géolocalise via Nominatim, les affiche
sur une carte OSMDroid avec tracé routier OSRM, et les enrichit de données
météo Open-Meteo. Ce module a été développé par \textbf{Anouar Soufyani}.

\section{Le formulaire de génération (TravelPathFragment)}

\ic{TravelPathFragment} présente un formulaire de génération suivi d'une liste
de parcours. Les deux parties coexistent dans le même fragment, séparées
visuellement. Le formulaire comprend sept champs :

\begin{tabularx}{\textwidth}{>{\bfseries}l X}
  \toprule
  Champ & Composant et comportement \\
  \midrule
  Destination
    & \ic{EditText}. Peut être pré-rempli depuis TravelShare via l'argument
      \ic{ARG\_CITY}. En cas de régénération, restauré depuis
      \ic{ARG\_REGEN\_ACTS}. \\
  Activités
    & Cinq \ic{CheckBox} : Culture, Restauration, Loisirs, Découverte,
      Shopping. Assemblées dans un \ic{LinkedHashSet} pour préserver l'ordre
      de sélection, qui détermine l'ordre des étapes. \\
  Budget max
    & \ic{SeekBar} de 10 à 300 €. Un \ic{OnSeekBarChangeListener} met à
      jour un \ic{TextView} en temps réel : "XX €". La valeur est lue avec
      \ic{Math.max(10, sbBudget.getProgress())}. \\
  Durée
    & \ic{SeekBar} de 1 à 12 h. Détermine le nombre d'étapes via \ic{maxStepsFor()} :
      2 h → 2 étapes, 4 h → 3, 6 h → 4, 8 h → 5 (ajusté ±1 selon l'effort). \\
  Effort
    & \ic{RadioGroup} : Facile (40 min/étape), Modéré (60 min), Intense (90 min).
      Influence aussi le nombre d'étapes (Facile = base−1, Intense = base+1). \\
  Lieux obligatoires
    & \ic{EditText} de saisie libre. Nominatim géocode chaque lieu et détecte
      son type OSM (\ic{class}/\ic{type}) pour estimer durée et coût adaptés. \\
  Tolérance météo
    & Trois \ic{CheckBox} : FROID, CHALEUR, HUMIDITE. Si au moins une est
      cochée (\ic{sensitiveWeather = true}), les loisirs proposés basculent
      vers des activités en intérieur. \\
  \bottomrule
\end{tabularx}

À l'appui sur "Générer", le bouton est désactivé, un Toast indique "Génération
en cours…", et \ic{viewModel.generatePlans()} est appelé. Le callback \ic{onDone}
réactive le bouton et bascule sur l'onglet "Mes parcours".

\screenshot{screen_travelpath_form.png}

\section{Affichage des parcours et onglets (TravelPathFragment)}

La liste des parcours est un \ic{RecyclerView} avec \ic{LinearLayoutManager}
alimenté par un \ic{PlanAdapter}. Chaque carte de parcours affiche : un bandeau
coloré selon le type (teal = économique, navy = équilibré, terracotta = confort),
le badge du type, la ville, le cœur like, les métriques (budget, durée, effort)
et les activités en format "Culture · Restauration · …".

Les deux onglets ("Mes parcours" / "Sauvegardés ❤️") sont implémentés via
\ic{MediatorLiveData}. Deux sources Room sont observées simultanément, mais
seule celle de l'onglet actif est propagée vers l'adapter.

\begin{lstlisting}
MediatorLiveData<List<TravelPlan>> merged = new MediatorLiveData<>();
LiveData<List<TravelPlan>> allPlans   = viewModel.getPlansForUser(userId);
LiveData<List<TravelPlan>> savedPlans = viewModel.getSavedPlansForUser(userId);

merged.addSource(allPlans,   list -> { if (!showSaved) merged.setValue(list); });
merged.addSource(savedPlans, list -> { if (showSaved)  merged.setValue(list); });
merged.observe(getViewLifecycleOwner(), adapter::setPlans);
\end{lstlisting}

En cliquant sur un onglet, \ic{showSaved} est inversé et
\ic{merged.setValue(source.getValue())} force le rafraîchissement immédiat
sans attendre un changement en base.

\screenshot[0.7\textwidth]{screen_travelpath_plans.png}

\section{Algorithme de génération (TravelPathViewModel)}

Toute la génération s'exécute sur le \ic{databaseWriteExecutor} pour ne pas
bloquer l'UI. Elle produit jusqu'à 3 objets \ic{TravelPlan} avec leurs
\ic{PlanStep} associés, composés de \textbf{vrais lieux existants} issus
d'OpenStreetMap.

\begin{lstlisting}
databaseWriteExecutor.execute(() -> {
    double[] cityCenter = geocode(city); // 1 appel Nominatim

    // 1. Overpass en parallele : 1 requete par activite
    //    Chaque requete retourne plusieurs lieux reels nommes
    //    Un lieu different est assigne a chaque type de plan
    ExecutorService pool = Executors.newFixedThreadPool(min(activities.size(), 5));
    for (String activity : activities) {
        pool.submit(() -> {
            // Recupere des lieux reels (restaurant, musee, parc...) via Overpass
            fetchRealPlaces(lat, lng, activity, sensitiveWeather);
            // Chaque type de plan (eco/equilibre/confort) obtient un lieu distinct
        });
    }
    pool.awaitTermination(25s);

    // 2. Construire et inserer les 3 plans
    for (String type : new String[]{"economique","equilibre","confort"}) {
        int maxSteps = maxStepsFor(durationHours, effort); // adapte au temps
        int stepDur  = stepDurationFor(effort);            // 40/60/90 min
        int budget   = (int)(budgetMax * ratio(type));     // 25/60/90 %
        List<PlanStep> steps = buildSteps(..., maxSteps, stepDur, budget, cache);
        steps = deduplicateSteps(steps); // supprime les doublons de lieux
        if (estimateBudget(steps) > budgetMax && !"economique".equals(type)) continue;
        geocodeStepsWithCenter(steps, city, cityCenter); // fallback si 0,0
        long planId = planDao.insertPlan(plan);
        for (PlanStep s : steps) stepDao.insertStep(s);
    }
    onDone.run();
});
\end{lstlisting}

\subsection{Lieux obligatoires : recherche par nom exact}

Quand l'utilisateur saisit un nom de lieu — restaurant, magasin, monument,
musée, ou tout autre établissement — dans le champ lieux obligatoires,
la méthode \ic{searchPlaceByName()} interroge Overpass avec un filtre regex
insensible à la casse (\ic{["name"{\raise.17ex\hbox{$\scriptstyle\sim$}}"...",i]}) dans un rayon de 16 km
autour du centre de la ville saisie.

Contrairement à un géocodage Nominatim classique qui retourne souvent
le centre-ville pour les enseignes, \textbf{Overpass retrouve la position GPS
exacte de l'établissement réel}, quelle que soit sa nature (restaurant,
boutique, musée, etc.). Ces coordonnées sont directement utilisées pour
placer le marqueur OSMDroid sur la carte — chaque lieu obligatoire
apparaît donc à son emplacement précis sur la carte, et deux lieux distincts
produisent deux pins à des emplacements différents.
Si Overpass échoue, Nominatim prend le relais.

\begin{lstlisting}
// Exemples de recherche par nom dans Paris (rayon 16 km)
node["name"~"McDonald.s",i](bbox); // → restaurant exact, lat/lon précis
node["name"~"Zara",i](bbox);       // → magasin exact
node["name"~"Louvre",i](bbox);     // → musée exact
// Chaque recherche retourne les coordonnées GPS du vrai établissement
\end{lstlisting}

\subsection{Adresses des étapes}

Chaque étape affiche son adresse réelle sous son nom. L'adresse est obtenue
de deux façons, dans cet ordre de priorité :

\begin{enumerate}
  \item \textbf{Tags OSM dans Overpass} : si le lieu a des tags \ic{addr:housenumber},
    \ic{addr:street}, \ic{addr:postcode} dans OSM, l'adresse est extraite
    directement depuis la réponse Overpass, sans appel réseau supplémentaire.
  \item \textbf{Reverse geocoding Nominatim} : si aucun tag d'adresse n'est
    disponible, la méthode statique \ic{reverseGeocode(lat, lng)} est appelée
    en arrière-plan depuis le \ic{databaseWriteExecutor} de l'adapter.
    L'adresse est persistée en base (\ic{planStepDao.updateAddress()}) pour
    fonctionner hors-ligne lors des consultations suivantes.
\end{enumerate}

\subsection{Déplacement sur la carte}

La carte OSMDroid est encapsulée dans un \ic{NestedScrollView}, qui interceptait
par défaut tous les événements tactiles, rendant impossible le déplacement
et le zoom. Un \ic{OnTouchListener} sur la \ic{MapView} appelle
\ic{requestDisallowInterceptTouchEvent(true)} dès qu'un doigt touche la carte,
puis relâche à \ic{ACTION\_UP}. La carte est ainsi pleinement interactive.

\subsection{Vrais lieux via Overpass API}

La méthode \ic{fetchRealPlaces()} interroge l'API Overpass (OpenStreetMap)
pour trouver des lieux réels nommés autour de la ville. Le filtre \ic{["name"]}
garantit que seuls les établissements avec un nom officiel sont retournés.

\begin{lstlisting}
// Exemple pour "Restauration" a Paris :
node["amenity"="restaurant"]["name"](48.74,2.20,48.97,2.46);
node["amenity"="cafe"]["name"](48.74,2.20,48.97,2.46);
// Retourne : "Le Procope", "Cafe de Flore", "Brasserie Lipp"...
\end{lstlisting}

Pour chaque activité, Overpass retourne plusieurs lieux nommés réels.
L'application en sélectionne un différent selon le type de plan —
ainsi le parcours économique, équilibré et confort n'affichent pas
le même établissement. Si Overpass ne répond pas, Nominatim prend le relais.

La méthode \ic{deduplicateSteps()} parcourt ensuite les étapes du plan et remplace
tout doublon (même nom) par le résultat alternatif disponible dans le cache.

\subsection{Budget, durée et effort réellement pris en compte}

Contrairement à une approche statique, les trois paramètres influencent
réellement la structure de chaque plan :

\begin{tabularx}{\textwidth}{>{\bfseries}l X}
  \toprule
  Paramètre & Impact concret \\
  \midrule
  Durée
    & \ic{maxStepsFor(durationHours, effort)} calcule le nombre d'étapes :
      de 2 (2 h Facile) à 6 (8 h Intense). \\
  Effort
    & \ic{stepDurationFor(effort)} fixe la durée de chaque étape :
      40 min (Facile), 60 min (Modéré), 90 min (Intense). \\
  Budget
    & Le coût de chaque étape est proportionnel au budget saisi :
      25 \% pour l'économique, 60 \% pour l'équilibré, 90 \% pour le confort. \\
  \bottomrule
\end{tabularx}

\subsection{Construction des étapes (buildSteps)}

La méthode \ic{buildSteps()} itère sur les activités sélectionnées et crée
une étape par activité, dans la limite de \ic{maxSteps}. Chaque étape est
d'abord remplie avec le vrai lieu OSM du cache Overpass, puis complétée
avec la durée et le coût calculés dynamiquement :

\begin{tabularx}{\textwidth}{>{\bfseries}l l X}
  \toprule
  Activité & Plan & Lieu réel trouvé + durée et coût estimés \\
  \midrule
  Culture & Économique & Monument ou musée gratuit — courte durée \\
  Culture & Équilibré  & Musée de la ville — durée normale \\
  Culture & Confort    & Galerie d'art ou musée premium — longue visite \\
  Restauration & Économique & Café ou restauration rapide \\
  Restauration & Équilibré  & Restaurant de la ville \\
  Restauration & Confort    & Restaurant (autre établissement) \\
  Loisirs & Météo sensible & Cinéma ou théâtre (activité couverte) \\
  Loisirs & Économique & Parc ou jardin public — gratuit \\
  Loisirs & Confort    & Spa ou centre de bien-être \\
  Découverte & Économique & Point de vue ou attraction gratuite \\
  Découverte & Confort   & Site touristique majeur \\
  Shopping & Économique & Marché local \\
  Shopping & Confort    & Boutique ou centre commercial \\
  \multicolumn{3}{l}{\small Durée et coût calculés dynamiquement selon le budget saisi et le niveau d'effort.}
  \bottomrule
\end{tabularx}

Si la durée dépasse 3 h et que "Restauration" n'est pas dans les activités
sélectionnées, une étape restauration est ajoutée automatiquement.

\section{Géocodage précis : Nominatim + viewbox + Haversine}

Nominatim retourne par défaut le résultat mondial le plus populaire. Sans
contrainte, "restaurant Paris" peut retourner un résultat à Paris, Texas.
La solution est un double filtre.

\textbf{Viewbox} : la ville est géocodée une fois en début de génération.
Ses coordonnées servent de centre pour construire un viewbox de ±0.4°
(≈44 km). Tous les appels suivants passent ce paramètre avec \ic{bounded=1}.

\textbf{Validation Haversine} : même avec le viewbox, un résultat peut
dépasser la zone. La formule de Haversine calcule la distance sphérique exacte
entre le résultat et le centre. Si elle dépasse 80 km, les coordonnées sont
rejetées (retour à 0, 0 — le point ne s'affiche pas sur la carte).

\begin{lstlisting}
private double[] geocodeNear(String query,
                              double centerLat, double centerLng) {
    double d = 0.4;
    String viewbox = String.format(Locale.US,
        "&viewbox=%.4f,%.4f,%.4f,%.4f&bounded=1",
        centerLng-d, centerLat-d, centerLng+d, centerLat+d);
    // Appel HTTP Nominatim + parse JSON manuel
    if (distanceKm(lat, lon, centerLat, centerLng) > 80)
        return new double[]{0, 0};
    return new double[]{lat, lon};
}

private double distanceKm(double lat1, double lng1,
                           double lat2, double lng2) {
    double dLat = Math.toRadians(lat2 - lat1);
    double dLng = Math.toRadians(lng2 - lng1);
    double a = Math.sin(dLat/2)*Math.sin(dLat/2)
             + Math.cos(Math.toRadians(lat1))
             * Math.cos(Math.toRadians(lat2))
             * Math.sin(dLng/2)*Math.sin(dLng/2);
    return 6371 * 2 * Math.atan2(Math.sqrt(a), Math.sqrt(1-a));
}
\end{lstlisting}

\begin{warnbox}[Délai inter-appels Nominatim]
Les CGU Nominatim imposent 1 req/s maximum. Un \ic{Thread.sleep(1100)} est
placé entre chaque appel de géocodage. Ce délai s'exécute sur le thread
background, l'interface reste réactive.
\end{warnbox}

\section{Affichage du parcours (PlanDetailFragment)}

\ic{PlanDetailFragment} est ouvert depuis le bouton "Voir le parcours" d'une
carte de plan. Il reçoit le \ic{planId} en argument et charge le plan et ses
étapes depuis Room.

\subsection{La carte OSMDroid et le tracé OSRM}

La carte est initialisée avec \ic{TileSourceFactory.MAPNIK} et le
multi-touch activé. Pour chaque étape valide (lat et lng non nuls), un
\ic{Marker} OSMDroid est créé avec titre et snippet, et ajouté aux overlays.

Le zoom est calculé manuellement car \ic{BoundingBox.fromGeoPoints()} n'existe
pas en OSMDroid 6.1.18. On parcourt tous les points pour calculer min/max des
latitudes et longitudes, puis on appelle \ic{mapView.zoomToBoundingBox()}.

Simultanément, si au moins 2 étapes sont valides et que la route n'a pas encore
été calculée, \ic{fetchOsrmRoute()} est lancé sur le \ic{databaseWriteExecutor}.

\begin{lstlisting}
private void fetchOsrmRoute(List<PlanStep> steps) {
    databaseWriteExecutor.execute(() -> {
        StringBuilder coords = new StringBuilder();
        for (PlanStep s : steps)
            coords.append(String.format(Locale.US,"%.6f,%.6f;",s.lng,s.lat));
        // URL OSRM : /route/v1/driving/...?overview=full&geometries=geojson
        // Reponse JSON -> routes[0].geometry.coordinates (GeoJSON)
        // GeoJSON : [lng, lat] -> GeoPoint(lat, lng)
        // Fallback si echec : lignes droites entre les etapes
        requireActivity().runOnUiThread(() -> drawRoutePolyline(routePoints));
    });
}

private void drawRoutePolyline(List<GeoPoint> points) {
    routeOverlay = new Polyline(mapView);
    routeOverlay.setPoints(points);
    routeOverlay.getOutlinePaint().setColor(0xFF1B3A5C);
    routeOverlay.getOutlinePaint().setStrokeWidth(9f);
    routeOverlay.getOutlinePaint().setStrokeCap(Paint.Cap.ROUND);
    mapView.getOverlays().add(0, routeOverlay); // Sous les marqueurs
    mapView.invalidate();
}
\end{lstlisting}

La polyline est stockée dans \ic{routeOverlay} comme champ du fragment. Ainsi,
si les étapes sont rechargées par LiveData (ce qui re-déclenche
\ic{addMarkersAndRoute()}), la route déjà calculée est réinsérée directement
sans refaire l'appel OSRM.

La réponse OSRM contient également \ic{routes[0].duration} (en secondes).
Ce temps est formaté (ex. "42 min" ou "1h05") et affiché dans un bloc
"TRAJET" dans la barre de métriques, rendu visible uniquement quand la route
est disponible.

\begin{lstlisting}
private String formatRouteDuration(double seconds) {
    int totalMin = (int) Math.round(seconds / 60.0);
    if (totalMin < 60) return totalMin + " min";
    int h = totalMin / 60, m = totalMin % 60;
    return m > 0 ? h + "h" + String.format("%02d", m) : h + "h";
}
\end{lstlisting}

\screenshot{screen_plan_detail.png}

\subsection{Liste des étapes (StepAdapter)}

Chaque étape est affichée avec une icône selon son type (🏛 Culture, 🍽
Restauration, 🎡 Loisirs, 🗺 Découverte, 🛍 Shopping, 📍 autre), le
créneau horaire, la durée, le nom, la description, le coût (ou "Gratuit"),
et les horaires d'ouverture typiques calculés statiquement selon le type.

\subsection{Mini-galerie TravelShare et vidéo dans les étapes}

Pour chaque étape, une requête \ic{photoDao.searchPhotosSync(keyword)} est
exécutée en background où \ic{keyword} est le type de l'étape (ex. "Culture").
Si des photos existent en base, un \ic{RecyclerView} horizontal avec Glide
les affiche. Si aucune photo ne correspond, un lien "Voir des photos" est
affiché. Ce lien ouvre \ic{ExplorerFragment} avec le nom de l'étape en
argument de recherche — c'est la passerelle TP$\to$TS.

Chaque étape dispose également d'un bouton "▶ Vidéo" qui ouvre YouTube via un
\ic{Intent.ACTION\_VIEW} avec l'URL de recherche construite dynamiquement
(\texttt{youtube.com/results?search\_query=NomÉtape+Type}). Cette approche
fournit du contenu vidéo pertinent sans aucune clé API ni hébergement.

\screenshot{screen_plan_detail.png}

\subsection{Mode hors-ligne}

Au chargement de \ic{TravelPathFragment}, \ic{ConnectivityManager} vérifie
l'état réseau. Si aucune connexion n'est active, une bannière terracotta
"📴 Hors-ligne — génération indisponible" s'affiche. Si l'utilisateur clique
sur "Générer" hors-ligne, un toast l'invite à consulter ses plans sauvegardés
et l'onglet "Sauvegardés ❤️" est activé automatiquement.

Les plans déjà générés restent entièrement consultables hors-ligne car ils sont
stockés dans Room. Seule la génération de nouveaux plans nécessite une
connexion (appels Nominatim pour le géocodage).

\subsection{La météo (Open-Meteo)}

Dès le chargement du détail, \ic{fetchWeather()} est lancé en background :
\begin{enumerate}
  \item Géocodage de \ic{plan.city} via Nominatim → lat/lng.
  \item Appel \ic{api.open-meteo.com/v1/forecast} avec \ic{current\_weather=true}.
  \item Parse du JSON : température, code météo WMO.
  \item Sur le thread UI : affichage dans \ic{layout\_weather} (icône,
    description, conseil, température).
\end{enumerate}

\begin{tabularx}{\textwidth}{l l X}
  \toprule
  Code WMO & Icône & Affichage \\
  \midrule
  0     & ☀  & Ciel dégagé — Parfait pour visiter \\
  1--3  & ⛅  & Partiellement nuageux — Bonne journée \\
  45--49 & -- & Brouillard — Visibilité réduite \\
  51--67 & -- & Pluie — Prévoyez un imperméable \\
  71--77 & -- & Neige — Habillez-vous chaudement \\
  80--82 & -- & Averses — Prévoyez un parapluie \\
  95+   & -- & Orage — Privilégiez les activités intérieures \\
  \bottomrule
\end{tabularx}

\subsection{Like / Sauvegarde}

Le bouton cœur dans le détail toggle \ic{plan.liked} et \ic{plan.saved}
(les deux sont synchronisés). \ic{viewModel.saveLikeAndSave(plan)} envoie
la mise à jour en base via \ic{planDao.updatePlan()}. Le plan lliké apparaît
alors dans l'onglet "Sauvegardés ❤️" de \ic{TravelPathFragment}.

\subsection{Export PDF}

Le bouton "Exporter PDF" déclenche sur le \ic{databaseWriteExecutor} la
construction d'un document \ic{PdfDocument} Android (API native, sans librairie).
Le PDF est rendu sur un \ic{Canvas} A4 (595×842 pt) avec trois objets \ic{Paint}
distincts (titre en bleu navy, sous-titres gris, étapes noir). Chaque étape
est inscrite avec créneau, nom, description, durée et coût. Le fichier est
sauvegardé dans \ic{Documents/TravelPath/parcours\_VILLE.pdf} et ouvert via
\ic{FileProvider} + \ic{Intent.ACTION\_VIEW}. Si aucune application PDF n'est
disponible, un Toast indique le chemin du fichier.

\subsection{Partage texte}

Le bouton "Partager" construit en background un résumé texte du parcours
(ville, toutes les étapes avec créneau et coût, budget total), puis lance
\ic{Intent.ACTION\_SEND} pour le partager via n'importe quelle application
(WhatsApp, email, SMS…).

\section{Régénération de parcours}

Le bouton "Régénérer" dans \ic{PlanDetailFragment} effectue deux opérations :
\begin{enumerate}
  \item \ic{getSupportFragmentManager().popBackStack()} pour revenir à la liste.
  \item Chargement de \ic{TravelPathFragment.newInstanceForRegen(plan)} qui
    reconstruit le formulaire avec tous les paramètres du plan courant.
\end{enumerate}

La restauration du formulaire parse la chaîne CSV \ic{plan.activities} pour
cocher les bonnes cases, lit \ic{plan.budgetEur} et \ic{plan.durationHours}
pour positionner les SeekBars, \ic{plan.effort} pour sélectionner le bon
RadioButton, et \ic{plan.requiredPlaces} + \ic{plan.weatherTolerances} pour
remplir les autres champs. Une bannière terracotta "🔄 Modifier et régénérer"
confirme visuellement la restauration.


% ==============================================================
\chapter{Passerelle — Intégration des deux modules}
% ==============================================================

\section{Principe}

La passerelle est ce qui différencie "Traveling" de deux applications séparées.
Elle permet une navigation naturelle entre TravelShare et TravelPath, avec partage
des données (lieux, photos) et une interface cohérente.

\section{TravelShare vers TravelPath}

Dans l'explorateur de TravelShare, un bouton "TravelPath" est visible en bas.
Si l'utilisateur cherchait "Montpellier" au moment de cliquer, cette destination
est automatiquement transmise au formulaire TravelPath via un argument Bundle.
Une bannière verte "📍 Destination importée depuis TravelShare" confirme l'import.

\begin{lstlisting}
view.findViewById(R.id.btn_explorer_travelpath).setOnClickListener(v ->
    requireActivity().getSupportFragmentManager()
        .beginTransaction()
        .replace(R.id.fragment_container, new TravelPathFragment())
        .addToBackStack(null)
        .commit()
);

String prefillCity = args != null ? args.getString(ARG_CITY, "") : "";
if (!prefillCity.isEmpty())
    ((EditText) view.findViewById(R.id.et_tp_city)).setText(prefillCity);
\end{lstlisting}

\section{TravelPath vers TravelShare}

Dans le détail d'un parcours TravelPath, chaque étape peut afficher des photos
TravelShare de la même catégorie. Le lien "Voir des photos" ouvre l'explorateur
TravelShare avec le nom de l'étape pré-rempli dans la barre de recherche.

\begin{lstlisting}
h.tvSeePhotos.setOnClickListener(v -> {
    ExplorerFragment explorer = new ExplorerFragment();
    Bundle args = new Bundle();
    args.putString(ExplorerFragment.ARG_SEARCH_QUERY, s.name);
    explorer.setArguments(args);
    fragment.requireActivity().getSupportFragmentManager()
        .beginTransaction()
        .replace(R.id.fragment_container, explorer)
        .addToBackStack(null)
        .commit();
});
\end{lstlisting}

\section{Régénération de parcours}

Un troisième flux de passerelle existe : depuis un parcours existant, l'utilisateur
peut régénérer un nouveau parcours en modifiant les paramètres. Tous les paramètres
(activités, budget, durée, effort, lieux) sont restaurés dans le formulaire.

\begin{lstlisting}
public static TravelPathFragment newInstanceForRegen(TravelPlan plan) {
    TravelPathFragment f = new TravelPathFragment();
    Bundle args = new Bundle();
    args.putString(ARG_CITY,           plan.city);
    args.putString(ARG_REGEN_ACTS,     plan.activities);
    args.putInt   (ARG_REGEN_BUDGET,   plan.budgetEur);
    args.putInt   (ARG_REGEN_DURATION, plan.durationHours);
    args.putString(ARG_REGEN_EFFORT,   plan.effort);
    args.putString(ARG_REGEN_PLACES,   plan.requiredPlaces);
    args.putString(ARG_REGEN_WEATHER,  plan.weatherTolerances);
    f.setArguments(args);
    return f;
}
\end{lstlisting}


% ==============================================================
\chapter{Modèle de données Room}
% ==============================================================

\section{Les 12 entités}

La base de données \ic{tp3\_database} (Room v19) contient 12 tables. Chaque table
correspond à une classe Java annotée \ic{@Entity}.

\begin{longtable}{>{\ttfamily\small}l >{\small}p{6.5cm} >{\centering\small}p{1.8cm}}
  \toprule
  \normalfont\textbf{Entité} & \textbf{Champs principaux et rôle} & \textbf{Module} \\
  \midrule
  \endhead
  User
    & id, login (unique), passwordHash, username,
      \textbf{avatarUri} (photo de profil), \textbf{bio} (biographie).
      Stocke les comptes utilisateurs.
    & Commun \\
  Photo
    & id, title, location, author, likes, lat, lng, date,
      category, tags, visibility (PUBLIC/GROUP), imageUri,
      groupId, \textbf{voiceUri}.
      Stocke toutes les photos publiées.
    & TS \\
  Comment
    & id, photoId, author, text, date.
      Commentaires attachés à une photo.
    & TS \\
  Group
    & id, name, description, creatorId.
      Groupes privés créés par les utilisateurs.
    & TS \\
  GroupMember
    & id, groupId, userId, userName, status (PENDING/MEMBER).
      Gère les membres et les demandes d'adhésion.
    & TS \\
  GroupMessage
    & id, groupId, senderId, senderName, content, date.
      Messages du chat de groupe.
    & TS \\
  Report
    & id, targetId, type, reason, reporterId, date.
      Signalements de contenu inapproprié.
    & TS \\
  NotificationPreference
    & id, userId, type (AUTEUR/LIEU/TAG/GROUPE), value.
      Préférences de notifications configurées par l'utilisateur.
    & TS \\
  AppNotification
    & id, targetUserId, type, message, groupId, isRead, date.
      Notifications in-app stockées pour consultation ultérieure.
    & TS \\
  Planning
    & id, userId, title, startDate, endDate.
      Plannings de voyage.
    & TP \\
  TravelPlan
    & id, userId, city, type (economique/equilibre/confort),
      activities (CSV), budgetEur, durationHours, effort,
      liked, saved, requiredPlaces, weatherTolerances, date.
      Parcours générés par TravelPath.
    & TP \\
  PlanStep
    & id, planId, stepOrder, name, type, timeSlot, durationMin,
      costEur, description, lat, lng, \textbf{openingHours}, \textbf{address}.
      Étapes individuelles d'un parcours, avec coordonnées GPS et adresse réelle.
    & TP \\
  \bottomrule
  \multicolumn{3}{l}{\small TS = TravelShare \quad TP = TravelPath}
\end{longtable}

\section{Relations entre les entités}

\begin{tcolorbox}[mybox,colback=codebg,colframe=bord,
  left=8pt,right=8pt,top=8pt,bottom=8pt]
{\ttfamily\footnotesize\color{navy!80}
\begin{verbatim}
User (1) ──────────────────────────────────────────────── (N)
  |                                                        |
  ├── (1:N) ──> Photo ──── (1:N) ──> Comment              |
  |              |                                         |
  |              └── (N:1) ──> Group ──── (1:N) ──> GroupMember
  |                               |                        |
  |                               └── (1:N) ──> GroupMessage
  |
  ├── (1:N) ──> AppNotification (isRead, type, groupId)
  ├── (1:N) ──> NotificationPreference (AUTEUR/LIEU/TAG)
  ├── (1:N) ──> Report
  ├── (1:N) ──> Planning
  └── (1:N) ──> TravelPlan ──── (1:N) ──> PlanStep (lat, lng, address)
\end{verbatim}}
\end{tcolorbox}

\section{Configuration de la base}

\begin{lstlisting}
@Database(
    entities = {User.class, Planning.class, Photo.class, Comment.class,
                Group.class, GroupMember.class, GroupMessage.class,
                Report.class, NotificationPreference.class,
                AppNotification.class, TravelPlan.class, PlanStep.class},
    version = 19,        // incremente a chaque changement de schema
    exportSchema = false
)
public abstract class AppDatabase extends RoomDatabase {
    // Pool de 4 threads : toutes les operations DB passent par ici
    public static final ExecutorService databaseWriteExecutor =
        Executors.newFixedThreadPool(4);

    // Singleton thread-safe (double-checked locking)
    public static AppDatabase getInstance(Context context) {
        if (INSTANCE == null) {
            synchronized (AppDatabase.class) {
                if (INSTANCE == null) {
                    INSTANCE = Room.databaseBuilder(context,
                                    AppDatabase.class, "tp3_database")
                        .fallbackToDestructiveMigration()
                        .build();
                }
            }
        }
        return INSTANCE;
    }
}
\end{lstlisting}


% ==============================================================
\chapter{Bilan et conclusion}
% ==============================================================

\section{Conclusion}

Ce projet nous a permis de construire une application Android complète, de bout en bout,
en partant d'une idée simple — faciliter la découverte et la planification de voyages —
jusqu'à une application fonctionnelle avec des données réelles.

Ce qui nous a le plus occupés, c'est la fiabilité. Obtenir de vrais lieux sur la carte,
des adresses précises, des itinéraires corrects, une synchronisation sans crash entre
threads background et UI — tout cela a demandé beaucoup d'itérations. Les difficultés
rencontrées (perte de données Room lors des migrations, géocodage erroné, événements
tactiles capturés par le ScrollView) ont toutes été résolues et ont renforcé notre
compréhension des mécanismes Android.

Nous avons tenu à n'utiliser que des services entièrement gratuits, ce qui nous a
poussés à explorer des alternatives open-source comme Overpass, Nominatim et OSRM —
des outils que nous n'aurions probablement pas découverts autrement.

\begin{okbox}[Ce qu'on retient]
Une application Android fonctionnelle avec deux modules complets et une passerelle,
des vrais lieux issus d'OpenStreetMap, un tracé routier réel, la météo du jour,
de l'annotation IA, de la pagination, du partage, un profil utilisateur — le tout
sans aucun service payant, développé en Java pur en suivant le pattern MVVM.
\end{okbox}


% ==============================================================
\chapter*{Annexes}
\addcontentsline{toc}{chapter}{Annexes}
% ==============================================================

\section*{A — Dépôt Git}

\begin{itemize}
  \item \textbf{URL} : \url{https://github.com/Soof2/Traveling.git}
  \item \textbf{Branche principale} : \ic{main}
\end{itemize}

\section*{B — APIs et ressources externes utilisées}

\begin{tabularx}{\textwidth}{l X l}
  \toprule
  \textbf{Service} & \textbf{Rôle} & \textbf{Coût} \\
  \midrule
  Nominatim OSM      & Géocodage d'adresses en coordonnées GPS
                     & Gratuit \\
  OSRM               & Calcul d'itinéraires routiers réels
                     & Gratuit \\
  Open-Meteo         & Données météo en temps réel
                     & Gratuit \\
  OSMDroid           & Librairie de cartographie Android
                     & Open-source \\
  Glide              & Chargement et mise en cache des images
                     & Open-source \\
  Overpass API (OSM) & Recherche de vrais lieux nommés (restaurants, musées…)
                     & Gratuit \\
  ML Kit (Google)    & Analyse IA on-device des photos (Image Labeling)
                     & Gratuit \\
  \bottomrule
\end{tabularx}

\section*{C — Permissions Android déclarées dans le Manifest}

\begin{itemize}
  \item \ic{android.permission.INTERNET}
  \item \ic{android.permission.ACCESS\_FINE\_LOCATION}
  \item \ic{android.permission.ACCESS\_COARSE\_LOCATION}
  \item \ic{android.permission.RECORD\_AUDIO}
  \item \ic{android.permission.CAMERA}
  \item \ic{android.permission.POST\_NOTIFICATIONS}
\end{itemize}

\section*{D — Environnement de développement}

\begin{tabularx}{\textwidth}{>{\bfseries}l X}
  \toprule
  Outil & Version \\
  \midrule
  Android Studio    & Hedgehog 2023.1+ \\
  SDK Android       & API 33 (target), API 24 (minimum) \\
  Java              & Java 8 (source compatibility) \\
  Room              & 2.6.x \\
  OSMDroid          & 6.1.18 \\
  Glide             & 4.16.x \\
  \bottomrule
\end{tabularx}

\end{document}
